% Section 4: From analytic to arbitrary slowly varying factors; proofs of Theorem 1 and Corollary 1.6 given Theorem 3.
% Labels defined here: sec:smoothing, sec:smoothing-gauss, sec:smoothing-fv, sec:smoothing-stab, sec:assembly,
%   lem:smoothing, lem:first-variation, lem:saturation, lem:apriori, prop:stability, rem:base-window,
%   eq:gauss-abs, eq:chaining, eq:fv-b, eq:fv-a, eq:diff-int, eq:gp, eq:head, eq:forms, eq:analytic-all-s.
\providecommand{\Ath}[1]{\textup{(A$_{#1}$)}}

\section{From analytic to arbitrary slowly varying factors}\label{sec:smoothing}

A slowly varying sequence is ``almost analytic''. Gaussian smoothing in the logarithmic variable turns $\log\ell$ into
an entire function $\Lambda_\ell$ whose derivatives tend to $0$ in every horizontal strip, and which differs from
$\log\ell$ by $o(1)$ at the integers (\S\ref{sec:smoothing-gauss}). The factor $\ell_1:=e^{\Lambda_\ell(\log\cdot)}$ is
therefore covered by the Riemann--Hilbert analysis, and it remains to compare the measures with factors $\ell$ and
$\ell_1$. Their ratio $\psi:=\ell/\ell_1$ tends to $1$. A multiplicative perturbation that tends to $1$ changes
$\log b_n^2$ by at most the supremum of $|\log\psi|$ over the bulk of the spectral weights, which sit at indices
$k\gtrsim n$ (\S\ref{sec:smoothing-fv}). The only danger is the head $k<n/M_0$, where $|\log\psi|$ need not be small;
there the spectral weights are exponentially small, because those atoms are saturated (\S\ref{sec:smoothing-stab}).

One cannot instead compare directly with the pure power $k^{-s}$: $\ell$ may be unbounded and need not be close to
$1$ anywhere, and controlling $\log\ell(k)-\log\ell(n)$ against the spectral weights far beyond $k\asymp n$ would require
information about the hard edge at the level of $1/n$ corrections. Comparing $\ell$ with its own smoothing avoids this.
\S\ref{sec:assembly} then proves Theorem~\ref{thm:main} and Corollary~\ref{thm:pure}, given Theorem~\ref{thm:engine}.

Throughout this section $p>0$ and $s>1$ are fixed, $x_k:=k^{-p}$, and for a positive sequence $w=(w_k)$ with
$\sum_kw_k<\infty$ we write $a_n(w)$, $b_n(w)$ for the Jacobi coefficients of $\sum_kw_k\delta_{x_k}$. With $p_n$ the
orthonormal polynomials (\S\ref{sec:prelim-op}), the spectral weights are
\[
  q_n(k):=w_k^{1/2}p_n(x_k),\qquad m_n(k):=q_n(k)^2,\qquad K_n(k):=\sum_{j<n}m_j(k)=w_k\mathsf K_n(x_k).
\]
Each $m_n$ is a probability vector on $\N$ (orthonormality of $p_n$), $0\le K_n(k)\le1$ by \eqref{eq:Kn-le-1}, and
$m_n=K_{n+1}-K_n$. For sequences $h_1,h_2$ on $\N$ we write $\langle h_1,h_2\rangle:=\sum_kh_1(k)h_2(k)$ whenever the sum converges
absolutely, and $x$ also denotes the sequence $(x_k)$.

\subsection{Gaussian smoothing}\label{sec:smoothing-gauss}

Let $L_\ell$, $G$, $\Lambda_\ell=L_\ell*G$ and $\rho_\ell$ be as in Definition~\ref{def:smoothing}; $G$ extends to an
entire function, and \Ath{\vartheta} is the hypothesis of Definition~\ref{def:A}.

Smoothing statements of this kind are familiar from the theory of regular variation; we include the short proof,
which uses only the sequence form of slow variation.

\begin{lemma}[Gaussian smoothing]\label{lem:smoothing}
Let $\ell$ be slowly varying and write $L=L_\ell$, $\Lambda=\Lambda_\ell$.
\begin{enumerate}
\item[\textup{(o)}] For every $\eta>0$, $|L(y)|\le A_\eta+\eta(|y|+\log2)$ for all $y\in\R$, with $A_\eta$ from
  Lemma~\ref{lem:potter}(ii). The function $\Lambda$ is entire and real on $\R$.
\item[\textup{(i)}] $\Lambda(\log k)-\log\ell(k)\to0$ as $k\to\infty$; that is, $\rho_\ell(m)\to0$ as $m\to\infty$.
\item[\textup{(ii)}] For every $\vartheta\ge0$, as $\Re\xi\to+\infty$ we have $\Lambda(\xi)-\Lambda(\Re\xi)\to0$ and
  $\Lambda^{(j)}(\xi)\to0$ for every $j\ge1$, uniformly in $|\Im\xi|\le\vartheta$; and $\Lambda$ is bounded on
  $\{\Re\xi\le0,\ |\Im\xi|\le\vartheta\}$.
\item[\textup{(iii)}] Consequently $\Lambda$ satisfies \Ath{\vartheta} for every $\vartheta>0$ (with any $c_0>0$). The sequence
  $\ell_1(k):=e^{\Lambda(\log k)}$ is positive, and $\psi:=\ell/\ell_1$ satisfies
  $\sup_{k\ge m}|\log\psi(k)|=\rho_\ell(m)\to0$.
\end{enumerate}
\end{lemma}

\begin{proof}
We prove (o), then (ii), from which (i) and (iii) follow.

\emph{(o).} Let $y\ge0$ and choose $k\ge1$ with $y\in[\log k,\log(k+1)]$. Then $L(y)$ is a convex combination of
$\log\ell(k)$ and $\log\ell(k+1)$, each of absolute value at most $A_\eta+\eta\log(k+1)\le A_\eta+\eta(y+\log2)$ by
Lemma~\ref{lem:potter}(ii). For $y<0$, $|L(y)|=|\log\ell(1)|\le A_\eta$. For $\xi=T+i\upsilon$ ($T,\upsilon\in\R$) and $y\in\R$,
\begin{equation}\label{eq:gauss-abs}
  |G(\xi-y)|=G(T-y)\,e^{\upsilon^2/2}.
\end{equation}
Since $L$ grows at most linearly, the integrand $L(y)G(\xi-y)$ is dominated, for $\xi$ in any bounded set, by an
integrable function of $y$; as it is entire in $\xi$ for each $y$, $\Lambda$ is entire, and its derivatives may be
taken under the integral sign. $\Lambda$ is real on $\R$ because $L$ and $G$ are.

\emph{Local oscillation of $L$.} For $T\ge T_0:=1+\log2$ put
\[
  \osc_L(T):=\sup\bigl\{|L(y)-L(y')|:\ y,y'\in[T-1,T+1]\bigr\},\qquad \osc_L^*(T):=\sup_{T'\ge T}\osc_L(T').
\]
Let $y\in[T-1,T+1]$ and $y\in[\log k,\log(k+1)]$. Then $k\ge e^{T-1}-1\ge e^{T-1}/2$ and $k+1\le e^{T+1}+1$. With
$m:=\lceil e^T\rceil$ and $\lambda:=2e^2$ we have $m/\lambda\le e^{T-2}\le e^{T-1}/2$ and $\lambda m\ge2e^{T+2}\ge
e^{T+1}+1$, so $k,k+1\in[m/\lambda,\lambda m]$. Hence $L(y)$ and $L(y')$ are convex combinations of values $\log\ell(k')$
with $k'\in[m/\lambda,\lambda m]$, and $\osc_L(T)\le2\osc_m(2e^2)$. Slow variation gives $\osc_L(T)\to0$, so
$\osc_L^*(T)$ is finite and $\osc_L^*(T)\to0$ as $T\to\infty$.

\emph{Chaining.} Let $T\ge2T_0$. If $y\ge T$, join $T$ to $y$ by $\lfloor y-T\rfloor$ steps of length $1$ and one shorter
step; every step lies in an interval $[T'-1,T'+1]$ with $T'\ge T$, so $|L(y)-L(T)|\le\osc_L^*(T)(1+|y-T|)$. If
$T/2\le y<T$, the same argument started at $y\ge T_0$ gives $|L(y)-L(T)|\le\osc_L^*(T/2)(1+|y-T|)$. If $y<T/2$, then
$|y|+T\le|y-T|+2T\le5|y-T|$, and (o) with $\eta=1$ gives $|L(y)-L(T)|\le2A_1+2\log2+|y|+T\le C_1(1+|y-T|)$ with
$C_1:=\max(2A_1+2,5)$. In all cases
\begin{equation}\label{eq:chaining}
  |L(y)-L(T)|\le\bigl(\osc_L^*(T/2)\,\mathbf 1_{\{y\ge T/2\}}+C_1\mathbf 1_{\{y<T/2\}}\bigr)(1+|y-T|).
\end{equation}

\emph{(ii).} The derivatives of $G$ are $G^{(j)}(w)=\mathsf h_j(w)G(w)$ with polynomials $\mathsf h_j$ of degree $j$.
Let $\bar{\mathsf h}_j$ be the polynomial whose coefficients are the absolute values of those of $\mathsf h_j$; then
$|\mathsf h_j(t+i\upsilon)|\le\bar{\mathsf h}_j(|t|+\vartheta)$ for $t\in\R$, $|\upsilon|\le\vartheta$. The function
$\xi\mapsto\int_\R G(\xi-y)\diff y$ is entire and equals $1$ on $\R$, hence everywhere; therefore
$\int_\R G^{(j)}(\xi-y)\diff y=\delta_{j0}$ for all $\xi\in\C$. For $\xi=T+i\upsilon$ with $|\upsilon|\le\vartheta$ and
$T\ge2T_0$, this, \eqref{eq:gauss-abs}, \eqref{eq:chaining} and the substitution $t=T-y$ give
\begin{align*}
  \bigl|\Lambda^{(j)}(\xi)-\delta_{j0}L(T)\bigr|&=\Bigl|\int_\R\bigl(L(y)-L(T)\bigr)G^{(j)}(\xi-y)\diff y\Bigr|\\
  &\le e^{\vartheta^2/2}\Bigl[\osc_L^*(T/2)\int_\R \mathsf g_j(t)\diff t+C_1\int_{|t|>T/2}\mathsf g_j(t)\diff t\Bigr],
\end{align*}
where $\mathsf g_j(t):=\bar{\mathsf h}_j(|t|+\vartheta)(1+|t|)G(t)$ is integrable. The right side tends to $0$ as $T\to\infty$, uniformly in
$|\upsilon|\le\vartheta$. For $j=0$, applying this at $\xi$ and at $T=\Re\xi$ and subtracting gives
$\sup_{|\Im\xi|\le\vartheta}|\Lambda(\xi)-\Lambda(\Re\xi)|\to0$. For $j\ge1$ it gives $\Lambda^{(j)}\to0$ uniformly in the
strip.

For the bound on $\Re\xi\le0$, note that $L(y)-L(0)=0$ for $y\le0$ and $|L(y)-L(0)|\le2A_1+2\log2+y$ for $y>0$. If
$T=\Re\xi\le0$ and $y>0$, then $|T-y|\ge y$, so $G(T-y)\le G(y)$. Hence
$|\Lambda(\xi)-L(0)|\le e^{\vartheta^2/2}\int_0^\infty(2A_1+2+y)G(y)\diff y$ for $|\Im\xi|\le\vartheta$.

\emph{(i).} Since $L(\log k)=\log\ell(k)$, the case $j=0$, $\Im\xi=0$ of the estimate in (ii) gives
$\Lambda(\log k)-\log\ell(k)\to0$.

\emph{(iii).} $\Lambda$ is entire, real on $\R$, bounded on $\{-c_0\le\Re\xi\le0,|\Im\xi|\le\vartheta\}$ by (ii), and
$\varepsilon_\Lambda(T)\to0$ by (ii) with $j=1$. So \Ath{\vartheta} holds. Finally
$\log\psi(k)=\log\ell(k)-\Lambda(\log k)$, so $\sup_{k\ge m}|\log\psi(k)|=\rho_\ell(m)$, which tends to $0$ by (i).
\end{proof}

\subsection{First variation}\label{sec:smoothing-fv}

The comparison of $\ell$ and $\ell_1$ runs along the path $w^{(\tau)}:=w^{(0)}\psi^\tau$, $\tau\in[0,1]$. The next lemma
computes the derivatives of the Jacobi coefficients along such a path. It is classical in substance (for $h=x$ it is
the Toda flow); we give the short proof.

\begin{lemma}[First variation]\label{lem:first-variation}
Let $x_k\in(0,1]$ be distinct, $v_k>0$ with $\sum_kv_k<\infty$, and let $h:\N\to\R$ be bounded. For $\tau\in\R$ put
$w^{(\tau)}_k:=v_ke^{\tau h(k)}$, and let $a_n,b_n,q_n,m_n$ denote the quantities of $w^{(\tau)}$. Then
$\tau\mapsto a_n$ ($n\ge0$) and $\tau\mapsto\log b_n^2$ ($n\ge1$) are $C^1$ on $\R$, and
\begin{align}
  \frac{\diff}{\diff\tau}\log b_n^2&=\langle h,m_n-m_{n-1}\rangle, \label{eq:fv-b}\\
  \frac{\diff}{\diff\tau}a_n&=\langle(x-a_n)h,m_n\rangle-2b_n\langle h,q_{n-1}q_n\rangle. \label{eq:fv-a}
\end{align}
\end{lemma}

\begin{proof}
Write $\sigma_\tau:=\sum_kw^{(\tau)}_k\delta_{x_k}$ and $\|h\|_\infty:=\sup_k|h(k)|$.

\emph{Smoothness.} The moments $\int x^i\diff\sigma_\tau=\sum_kv_ke^{\tau h(k)}x_k^i$ are $C^\infty$ in $\tau$, with
derivatives computed term by term: the $r$-th termwise derivative is bounded by $v_k\|h\|_\infty^re^{|\tau|\|h\|_\infty}$, which is summable
uniformly for $\tau$ in compact sets. The coefficients of $P_n$ solve the linear system $\int P_nx^i\diff\sigma_\tau=0$
($i<n$), whose determinant is the Hankel determinant $D_n(\sigma_\tau)>0$; by Cramer's rule they are $C^\infty$ in
$\tau$. So are $h_n=\int P_n^2\diff\sigma_\tau>0$, $p_n=P_n/\sqrt{h_n}$, $a_n$ and $b_n^2=h_n/h_{n-1}$. Moreover, if
$g(x;\tau)$ is a polynomial in $x$ whose coefficients are $C^1$ in $\tau$, then
\begin{equation}\label{eq:diff-int}
  \frac{\diff}{\diff\tau}\int g\diff\sigma_\tau=\int\partial_\tau g\diff\sigma_\tau+\sum_kh(k)w^{(\tau)}_kg(x_k;\tau),
\end{equation}
since $\int g\diff\sigma_\tau$ is a finite linear combination of moments.

\emph{The $b$-formula.} Apply \eqref{eq:diff-int} to $g=P_n^2$. Since $P_n$ is monic, $\partial_\tau P_n$ has degree
less than $n$ and is orthogonal to $P_n$. Hence $\frac{\diff}{\diff\tau}h_n=\sum_kh(k)w^{(\tau)}_kP_n(x_k)^2=h_n\langle
h,m_n\rangle$, i.e.\ $\frac{\diff}{\diff\tau}\log h_n=\langle h,m_n\rangle$. Since $\log b_n^2=\log h_n-\log h_{n-1}$, this
is \eqref{eq:fv-b}.

\emph{The $a$-formula.} Write $\dot p_n:=\partial_\tau p_n$, a polynomial of degree at most $n$, and expand
$\dot p_n=\sum_{j\le n}\mathsf c_jp_j$ with $\mathsf c_j=\int\dot p_np_j\diff\sigma_\tau$. Differentiating
$\int p_np_j\diff\sigma_\tau=\delta_{nj}$ with \eqref{eq:diff-int} gives
$\int\dot p_np_j\diff\sigma_\tau+\int p_n\dot p_j\diff\sigma_\tau+\langle h,q_nq_j\rangle=0$. For $j<n$ the middle term
vanishes ($\deg\dot p_j\le j<n$), so $\mathsf c_j=-\langle h,q_nq_j\rangle$; for $j=n$ we get $\mathsf c_n=-\frac12\langle
h,m_n\rangle$. Now $a_n=\int xp_n^2\diff\sigma_\tau$ by \eqref{eq:coef-int}, and by \eqref{eq:diff-int} and
\eqref{eq:coef-int},
\begin{align*}
  \frac{\diff a_n}{\diff\tau}&=2\int xp_n\dot p_n\diff\sigma_\tau+\langle xh,m_n\rangle
  =2\mathsf c_na_n+2\mathsf c_{n-1}b_n+\langle xh,m_n\rangle\\
  &=-a_n\langle h,m_n\rangle-2b_n\langle h,q_{n-1}q_n\rangle+\langle xh,m_n\rangle,
\end{align*}
which is \eqref{eq:fv-a}. (For $n=0$ the term with $\mathsf c_{n-1}$ is absent, consistent with $b_0=0$.) All the sums
converge absolutely, since $h$ and $x$ are bounded, $m_n$ is a probability vector and $|q_{n-1}q_n|\le\frac12(m_{n-1}+m_n)$.
\end{proof}

\subsection{\texorpdfstring{The head is saturated; stability under $\ell/\ell_1\to1$}{The head is saturated; stability}}\label{sec:smoothing-stab}

For $p>0$ and $0<y<1$ let
\begin{equation}\label{eq:gp}
  g_p(y):=(p-y)\log\frac1y+(1-y)\log(1-y)-p+(p-1)y,\qquad
  g_p^+:=\begin{cases}g_p,&p\ge1,\\ g_p-\log\frac1p,&0<p<1.\end{cases}
\end{equation}
For $y<p/2$ the first term is at least $\frac p2\log\frac1y$, and the other terms are bounded; so
$g_p^+(y)\to+\infty$ as $y\to0^+$. Hence there is $y_1>0$ with $g_p^+\ge1$ on $(0,y_1]$. We fix the \emph{head margin} $M_0=M_0(p)$ used in
the rate \eqref{eq:main-rate}: it is the least integer $M_0\ge4$ such that $g_p^+\ge1$ on $(0,2/M_0]$.

The next lemma says that the atoms $x_k$ with $k$ much smaller than $n$ are almost fully occupied by the first $n$
orthonormal polynomials. Its proof combines Christoffel's extremal property with a Lagrange test polynomial.

\begin{lemma}[Saturation of the head]\label{lem:saturation}
Let $w_k>0$ with $\|w\|_1:=\sum_kw_k<\infty$, and consider $\sum_kw_k\delta_{k^{-p}}$. For integers $1\le k<n$ and $y:=k/n$,
\[
  1-K_n(k)\le\frac{\|w\|_1}{w_k}\Bigl(\frac{n+1}k\Bigr)^{2(1-p)_+}e^{-2ng_p^+(y)} .
\]
Moreover $m_n(k)\le1-K_n(k)$ for all $k$ and $n$.
\end{lemma}

\begin{proof}
The last statement holds because $m_n(k)=K_{n+1}(k)-K_n(k)$ and $K_{n+1}(k)\le1$. For the main bound we use the test
polynomial
\[
  P(x):=\prod_{j\le n,\,j\ne k}\frac{x-x_j}{x_k-x_j},
\]
which has degree $n-1$, satisfies $P(x_k)=1$, and vanishes at $x_j$ for $j\le n$, $j\ne k$. For $i>n$ and $j\le n$ we have
$0<x_i<x_j$, so $|x_i-x_j|\le x_j$ and $|P(x_i)|\le\mathcal P:=\prod_{j\le n,\,j\ne k}x_j/|x_k-x_j|$. Therefore
$\int P^2\diff\sigma\le w_k+\|w\|_1\mathcal P^2$, and \eqref{eq:christoffel} gives $1/\mathsf K_n(x_k)\le w_k+\|w\|_1\mathcal P^2$. Hence
$K_n(k)=w_k\mathsf K_n(x_k)\ge w_k/(w_k+\|w\|_1\mathcal P^2)$ and
\[
  1-K_n(k)\le\frac{\|w\|_1\mathcal P^2}{w_k+\|w\|_1\mathcal P^2}\le\frac{\|w\|_1}{w_k}\mathcal P^2 .
\]

It remains to bound $\mathcal P$. Here $x_j/|x_k-x_j|=k^p/|k^p-j^p|$. Put $\varsigma_p:=\min(1,p)$. We claim
$|k^p-j^p|\ge \varsigma_p\max(j,k)^{p-1}|k-j|$. For $p\ge1$ and $j<k$: $k^p-j^p=k^{p-1}k-j^{p-1}j\ge k^{p-1}(k-j)$; for $p\ge1$
and $j>k$ the same computation gives $j^p-k^p\ge j^{p-1}(j-k)$. For $p<1$ the mean value theorem gives
$|k^p-j^p|=p\xi^{p-1}|k-j|$ with $\xi$ between $j$ and $k$, and $\xi^{p-1}\ge\max(j,k)^{p-1}$. Hence
\[
  \mathcal P\le \varsigma_p^{-(n-1)}\prod_{j<k}\frac k{k-j}\prod_{j=k+1}^n\frac{k^p}{j^{p-1}(j-k)}
  =\varsigma_p^{-(n-1)}\,\frac{k^{k-1}}{(k-1)!}\cdot\frac{k^{p(n-k)}\,(n!/k!)^{1-p}}{(n-k)!} .
\]
We bound the three factors. First, $k^{k-1}/(k-1)!=k^k/k!\le e^k$. Second, $\log(n-k)!\ge(n-k)\log(n-k)-(n-k)$.
Third, $S:=\log(n!/k!)=\sum_{j=k+1}^n\log j$ satisfies
\[
  I_{k,n}:=\int_k^n\log t\diff t\le S\le I_{k,n}+\log\frac{n+1}k,
\]
because $\log$ is increasing ($\log j\ge\int_{j-1}^j\log$ and $\log j\le\int_j^{j+1}\log$). We use the lower bound if
$p\ge1$ (then $1-p\le0$) and the upper bound if $p<1$; in both cases $(1-p)S\le(1-p)I_{k,n}+(1-p)_+\log\frac{n+1}k$.
With $I_{k,n}=n\log n-n-k\log k+k$ and $k=yn$, the terms in $\log n$ cancel and
\[
  k+p(n-k)\log k+(1-p)I_{k,n}-(n-k)\log(n-k)+(n-k)=-n\,g_p(y),
\]
as a direct expansion shows. (The coefficient of $n\log n$ is $p(1-y)+(1-p)(1-y)-(1-y)=0$.) Hence
\[
  \log\mathcal P\le-n\,g_p(y)+(n-1)\log\frac1{\varsigma_p}+(1-p)_+\log\frac{n+1}k .
\]
For $p\ge1$, $\varsigma_p=1$ and squaring gives the claim. For $p<1$, $\varsigma_p=p$ and $(n-1)\log\frac1p\le n\log\frac1p$, so
$\mathcal P^2\le e^{-2n(g_p(y)-\log(1/p))}((n+1)/k)^{2(1-p)}$, which is the claim with $g_p^+$.
\end{proof}

We now fix the setting of the comparison. Let $\ell$ be slowly varying, let $\Lambda_\ell$, $\ell_1$, $\psi=\ell/\ell_1$
be as in Lemma~\ref{lem:smoothing}, and put
\[
  \bar\rho:=\sup_k|\log\psi(k)|=\rho_\ell(1)<\infty,\qquad
  w^{(\tau)}_k:=k^{-s}\ell_1(k)\psi(k)^\tau\quad(\tau\in[0,1]).
\]
So $w^{(0)}$ is the weight of $\sigma_{p,s,\ell_1}$ and $w^{(1)}$ that of $\sigma_{p,s,\ell}$. Since
$\ell_1\psi^\tau=\ell\psi^{\tau-1}$, we have $e^{-\bar\rho}\ell(k)\le\ell_1(k)\psi(k)^\tau\le e^{\bar\rho}\ell(k)$, so every
$w^{(\tau)}$ has finite total mass. We call the indices $k<n/M_0$ the \emph{head}, and let
\[
  \mathfrak h_n:=\sup_{\tau\in[0,1]}\ \max_{m\in\{n-1,n\}}\ \sum_{k<n/M_0}m^{(\tau)}_m(k)
\]
be the largest head mass of the spectral weights $m^{(\tau)}_{n-1}$, $m^{(\tau)}_n$ of $w^{(\tau)}$.

\begin{lemma}[A priori size]\label{lem:apriori}
In this setting there are constants $C$ and $n_0$, depending on $p,s,\ell$, such that for all $n\ge2$
\begin{equation}\label{eq:head}
  \mathfrak h_n\le Cn^{s+5}e^{-2n},
\end{equation}
and for all $n\ge n_0$ and $\tau\in[0,1]$,
\[
  a_n(w^{(\tau)})\le2M_0^pn^{-p},\qquad b_n(w^{(\tau)})\le2M_0^pn^{-p}.
\]
\end{lemma}

\begin{proof}
\emph{Head.} By Lemma~\ref{lem:potter}(ii) with $\eta=1$, we have
$\|w^{(\tau)}\|_1\le e^{\bar\rho}\sigma_{p,s,\ell}(\R)<\infty$ and
$w^{(\tau)}_k\ge e^{-\bar\rho}k^{-s}\ell(k)\ge e^{-\bar\rho-A_1}k^{-s-1}$, so $\|w^{(\tau)}\|_1/w^{(\tau)}_k\le C'k^{s+1}$ with
$C':=e^{2\bar\rho+A_1}\sigma_{p,s,\ell}(\R)$. Let $n\ge2$, $m\in\{n-1,n\}$ and $k<n/M_0$. Then $k<n/4<m$, and $k/m\le
n/(M_0(n-1))\le2/M_0$, so $g_p^+(k/m)\ge1$. Lemma~\ref{lem:saturation}, with $(m+1)/k\le n+1$ and $2(1-p)_+\le2$, gives
\[
  m^{(\tau)}_m(k)\le1-K^{(\tau)}_m(k)\le C'k^{s+1}(n+1)^2e^{-2(n-1)} .
\]
Summing over the fewer than $n$ indices $k<n/M_0$ gives \eqref{eq:head}.

\emph{Size.} By \eqref{eq:coef-int}, $a_m=\langle x,m_m\rangle$ and $b_n=\langle xq_{n-1},q_n\rangle\le\frac12\langle
x,m_{n-1}+m_n\rangle$, since $x\ge0$ and $|q_{n-1}q_n|\le\frac12(m_{n-1}+m_n)$. For $k\ge n/M_0$ we have
$x_k=k^{-p}\le(M_0/n)^p$, and $x_k\le1$ always. Hence, for $m\in\{n-1,n\}$,
$\langle x,m^{(\tau)}_m\rangle\le(M_0/n)^p+\mathfrak h_n$, and both $a_n$ and $b_n$ are at most $(M_0/n)^p+\mathfrak h_n$. By
\eqref{eq:head}, $\mathfrak h_n\le(M_0/n)^p$ for $n\ge n_0$.
\end{proof}

\begin{proposition}[Stability under $\ell/\ell_1\to1$]\label{prop:stability}
Let $p>0$, $s>1$, let $\ell$ be slowly varying, and let $\ell_1:=e^{\Lambda_\ell(\log\cdot)}$. There are constants $C$ and
$n_0$, depending on $p,s,\ell$, such that for $n\ge n_0$
\begin{align*}
  \bigl|\log b_n^2(\sigma_{p,s,\ell})-\log b_n^2(\sigma_{p,s,\ell_1})\bigr|
    &\le2\rho_\ell(\lceil n/M_0\rceil)+Cn^{s+5}e^{-2n},\\
  n^p\bigl|a_n(\sigma_{p,s,\ell})-a_n(\sigma_{p,s,\ell_1})\bigr|
    &\le6M_0^p\,\rho_\ell(\lceil n/M_0\rceil)+Cn^{p+s+5}e^{-2n}.
\end{align*}
\end{proposition}

\begin{proof}
Apply Lemma~\ref{lem:first-variation} with $v_k=k^{-s}\ell_1(k)$ and $h=\log\psi$, which is bounded by $\bar\rho$; then
$w^{(\tau)}$ is the path defined above. Write $\rho:=\rho_\ell(\lceil n/M_0\rceil)$. For $k\ge n/M_0$ we have
$k\ge\lceil n/M_0\rceil$, so $|h(k)|\le\rho$; for $k<n/M_0$, $|h(k)|\le\bar\rho$.

\emph{The $b$-part.} By \eqref{eq:fv-b} and the fundamental theorem of calculus,
\[
  \bigl|\log b_n^2(w^{(1)})-\log b_n^2(w^{(0)})\bigr|\le\sup_{\tau\in[0,1]}\sum_k|h(k)|\bigl(m^{(\tau)}_n(k)+m^{(\tau)}_{n-1}(k)\bigr)
  \le2\rho+2\bar\rho\mathfrak h_n,
\]
because $m^{(\tau)}_n$ and $m^{(\tau)}_{n-1}$ are probability vectors. Lemma~\ref{lem:apriori} bounds $\mathfrak h_n$.

\emph{The $a$-part.} By \eqref{eq:fv-a} and $2|q_{n-1}q_n|\le m_{n-1}+m_n$,
\[
  \Bigl|\frac{\diff a_n}{\diff\tau}\Bigr|\le\sum_k|h(k)|\Bigl(|x_k-a_n|\,m_n(k)+b_n\bigl(m_{n-1}(k)+m_n(k)\bigr)\Bigr).
\]
Let $n\ge n_0$, where $n_0$ is as in Lemma~\ref{lem:apriori}, enlarged so that $2M_0^pn_0^{-p}\le1$. For $k\ge n/M_0$, $0\le x_k\le(M_0/n)^p$ and $0\le a_n\le2M_0^pn^{-p}$,
so $|x_k-a_n|\le2M_0^pn^{-p}$; also $b_n\le2M_0^pn^{-p}$. These $k$ contribute at most
$\rho\bigl(2M_0^pn^{-p}+2\cdot2M_0^pn^{-p}\bigr)=6M_0^p\rho\,n^{-p}$. For $k<n/M_0$ use $|x_k-a_n|\le1$ and
$b_n\le1$; these contribute at most $3\bar\rho\mathfrak h_n$. Integrating over $\tau\in[0,1]$ and
multiplying by $n^p$ gives the second bound, again with Lemma~\ref{lem:apriori}.
\end{proof}


\subsection{\texorpdfstring{Proof of Theorem~\ref{thm:main}, assuming Theorem~\ref{thm:engine}}{Proof of Theorem 1, assuming Theorem 3}}\label{sec:assembly}

For a sequence of measures with Jacobi coefficients $a_n,b_n$ and a constant $c>0$, the two forms of the rate
statement are equivalent:
\begin{equation}\label{eq:forms}
  \Bigl\{n^pb_n=c\bigl(1+O(1/n)\bigr),\ \ \frac{a_n}{2b_n}=1+O(1/n)\Bigr\}
  \iff |n^pa_n-2c|+|n^pb_n-c|=O(1/n).
\end{equation}
Indeed, from the left side, $n^pa_n=2n^pb_n\cdot a_n/(2b_n)=2c(1+O(1/n))^2$; from the right side,
$n^pb_n=c+O(1/n)$ and $a_n/(2b_n)=(2c+O(1/n))/(2c+O(1/n))$. We write $c=c(p)$.

\emph{Step 1: analytic factors, every $s>1$.} Let $\ell_1$ satisfy \Ath{\vartheta} for some $\vartheta>0$ and let $s>1$. We claim
\begin{equation}\label{eq:analytic-all-s}
  |n^pa_n(\sigma_{p,s,\ell_1})-2c|+|n^pb_n(\sigma_{p,s,\ell_1})-c|=O(1/n).
\end{equation}
Let $m\ge0$ be the least integer with $s_0:=s-mp<1+2p$. If $m=0$ then $s_0=s>1$; if $m\ge1$ then $s-(m-1)p\ge1+2p$, so
$s_0\ge1+p$. In both cases $s_0\in(1,1+2p)$. Theorem~\ref{thm:engine} at $s_0$ and \eqref{eq:forms} give
\eqref{eq:analytic-all-s} for $\sigma_{p,s_0,\ell_1}$. By \eqref{eq:shift-identity},
$\sigma_{p,s,\ell_1}=x^m\sigma_{p,s_0,\ell_1}$, and each $x^j\sigma_{p,s_0,\ell_1}=\sigma_{p,s_0+jp,\ell_1}$ is a positive
measure on $(0,1]$ with infinite support and finite moments. Theorem~\ref{thm:shift}, applied $m$ times with the regular
rate $\eta_n=1/n$, gives \eqref{eq:analytic-all-s} for $\sigma_{p,s,\ell_1}$.

\emph{Step 2: Theorem~\ref{thm:main}(ii).} If $\ell$ itself satisfies \Ath{\vartheta} for some $\vartheta>0$, Step 1 with $\ell_1=\ell$ and
\eqref{eq:forms} give $n^pb_n=c(1+O(1/n))$ and $a_n/(2b_n)=1+O(1/n)$.

\emph{Step 3: Corollary~\ref{thm:pure}.} The constant sequence $\ell\equiv1$ is $e^{\Lambda(\log k)}$ with $\Lambda\equiv0$,
which satisfies \Ath{\vartheta} for every $\vartheta$ (with $\varepsilon_\Lambda\equiv0$). Step 2 applies.

\emph{Step 4: the rate \eqref{eq:main-rate}.} Let $\ell$ be slowly varying and $\ell_1:=e^{\Lambda_\ell(\log\cdot)}$. By
Lemma~\ref{lem:smoothing}(iii), $\Lambda_\ell$ satisfies \Ath{\vartheta} for every $\vartheta>0$, so Step 1 gives
\eqref{eq:analytic-all-s} for $\sigma_{p,s,\ell_1}$. Write $\rho:=\rho_\ell(\lceil n/M_0\rceil)$ and
$\Delta_n:=\log b_n^2(\sigma_{p,s,\ell})-\log b_n^2(\sigma_{p,s,\ell_1})$. By Proposition~\ref{prop:stability},
$|\Delta_n|\le2\rho+Cn^{s+5}e^{-2n}$ for $n\ge n_0$; in particular $|\Delta_n|$ is bounded, since $\rho\le\rho_\ell(1)$.
Hence $|e^{\Delta_n/2}-1|\le C'|\Delta_n|$ and
\[
  |n^pb_n(\sigma_{p,s,\ell})-c|\le|n^pb_n(\sigma_{p,s,\ell_1})-c|+n^pb_n(\sigma_{p,s,\ell_1})\,|e^{\Delta_n/2}-1|
  =O(1/n)+O(\rho)+O(n^{s+5}e^{-2n}).
\]
Similarly, by the second bound of Proposition~\ref{prop:stability},
$|n^pa_n(\sigma_{p,s,\ell})-2c|\le|n^pa_n(\sigma_{p,s,\ell_1})-2c|+6M_0^p\rho+Cn^{p+s+5}e^{-2n}$. The exponentially small
terms are $O(1/n)$, and finitely many $n$ only affect the constant. This is
\[
  |n^pb_n(\sigma_{p,s,\ell})-c|+|n^pa_n(\sigma_{p,s,\ell})-2c|=O\bigl(1/n+\rho_\ell(\lceil n/M_0\rceil)\bigr).
\]

\emph{Step 5: Theorem~\ref{thm:main}(i)} follows from \eqref{eq:main-rate}, because $\rho_\ell(m)\to0$ by Lemma~\ref{lem:smoothing}(i).
\qed

\medskip
\emph{On the order of the steps.} The smoothing is applied last, at the final exponent $s$, after the shift. This is
possible because Proposition~\ref{prop:stability} holds for every $p>0$ and $s>1$, while Theorem~\ref{thm:shift} is
applied only to the analytic factor $\ell_1$, with the rate $1/n$ of Theorem~\ref{thm:engine}. Smoothing first and
shifting afterwards would also work, but the shift would then have to carry the rate $1/n+\rho_\ell(\lceil n/M_0\rceil)$,
which in general is not regular and must be replaced by a regular majorant (Remark~\ref{rem:majorant}).

\begin{remark}[Which part of Theorem~\ref{thm:engine} is used]\label{rem:base-window}
Any half-open $s$-window of length $p$ inside $(1,1+2p)$ can serve as the base in Step 1. For $s>1+p/2$ one may take
$s_0$ to be the unique point of $s-p\N_0$ in $(1+p/2,1+3p/2]$, so that $\alpha=(s_0-1-p)/p\in(-\frac12,\frac12]$. Thus for
every $s>1+p/2$ only the case $|\alpha|\le\frac12$ of Theorem~\ref{thm:engine} is used, and the range $\frac12<|\alpha|<1$
is needed only for $1<s\le1+p/2$, where no shift takes place. The case $\alpha=0$ cannot be avoided: for
$s\in1+p+p\N_0$ every admissible base point has $s_0\in1+p\mathbb Z$, i.e.\ $\alpha\in\mathbb Z\cap(-1,1)=\{0\}$.
\end{remark}
